\section{An integral identity for the hypergeometric function}
\label{m:app:hyp}

The route to \cref{m:eq:Psidef} taken in \cref{m:app:absbirthdeath} expands the
characteristic $x(\tau)$ in a geometric series and integrates term by term. An
alternative is to integrate the quotient \cref{m:eq:xabd} directly, splitting the
numerator into two pieces, each of the form
\begin{equation}
\label{m:eq:Iform}
I(\tau) = \int\frac{\ee^{h\tau}}{P+Q\,\ee^{k\tau}}\,\dd\tau
\end{equation}
with $h,k,P,Q$ constant. That route is longer and easier to get wrong, but the
identity it rests on is useful in its own right.

\begin{proposition}
\label{m:prop:hypIden}
For $h/k>0$ and $-\tfrac{Q}{P}\ee^{k\tau}\notin[1,\infty)$,
\begin{equation}
\label{m:eq:hypIden}
\int\frac{\ee^{h\tau}}{P+Q\ee^{k\tau}}\,\dd\tau
= \frac{\ee^{h\tau}}{Ph}\,{}_2F_1\lrb{1,\frac{h}{k};\frac{h}{k}+1;\,-\frac{Q}{P}\ee^{k\tau}}
+ \text{const} .
\end{equation}
\end{proposition}

\begin{proof}
We use the Euler integral representation of the hypergeometric function
\cite{Olver2010},
\begin{equation}
\label{m:eq:idenHyp}
{}_2F_1(a,b;c;z)
= \frac{\Gamma(c)}{\Gamma(b)\Gamma(c-b)}\int_0^1 u^{b-1}(1-u)^{c-b-1}(1-zu)^{-a}\,\dd u ,
\end{equation}
$\Gamma$ being the gamma function.

Start from \cref{m:eq:Iform} and substitute $v = \ee^{k\tau}$, so that
$\dd\tau = k^{-1}v^{-1}\dd v$:
\begin{equation*}
I = \frac1k\int\frac{v^{h/k}}{P+Qv}\,v^{-1}\dd v
= \frac{1}{Pk}\int\frac{v^{h/k-1}}{1+\frac{Q}{P}v}\,\dd v .
\end{equation*}
Fix the antiderivative by writing this as a definite integral from $0$ to $v$ in a
dummy variable $w$; the integral converges at the lower limit precisely because
$h/k>0$:
\begin{equation*}
I = \frac{1}{Pk}\int_0^v\frac{w^{h/k-1}}{1+\frac{Q}{P}w}\,\dd w .
\end{equation*}
Now rescale, $w = vu$, so that $\dd w = v\,\dd u$ with $v$ held fixed:
\begin{equation}
\label{m:eq:backat}
I = \frac{1}{Pk}\int_0^1\frac{(vu)^{h/k-1}}{1+\frac{Q}{P}vu}\,v\,\dd u
= \frac{v^{h/k}}{Pk}\int_0^1\frac{u^{h/k-1}}{1+\frac{Q}{P}vu}\,\dd u .
\end{equation}
The remaining integrand matches that of \cref{m:eq:idenHyp} with $b = h/k$,
$c = h/k+1$, $a=1$ and $z = -\tfrac{Q}{P}v$, since then $(1-u)^{c-b-1} = 1$ and
$(1-zu)^{-a} = \bigl(1+\tfrac{Q}{P}vu\bigr)^{-1}$. Hence, using
$\Gamma(z+1) = z\Gamma(z)$,
\begin{equation*}
\int_0^1\frac{u^{h/k-1}}{1+\frac{Q}{P}vu}\,\dd u
= \frac{\Gamma\!\lrb{\frac hk}\Gamma(1)}{\Gamma\!\lrb{\frac hk+1}}\,
{}_2F_1\lrb{1,\frac hk;\frac hk+1;-\frac QP v}
= \frac kh\,{}_2F_1\lrb{1,\frac hk;\frac hk+1;-\frac QP v}.
\end{equation*}
Substituting into \cref{m:eq:backat} gives
\begin{equation*}
I = \frac{v^{h/k}}{Ph}\,{}_2F_1\lrb{1,\frac hk;\frac hk+1;-\frac QP v},
\end{equation*}
and restoring $v = \ee^{k\tau}$ gives \cref{m:eq:hypIden}.
\end{proof}

\begin{remark}
Applied to the two pieces of \cref{m:eq:xabd} --- the first with $h=\alpha$, the
second with $h = \alpha+b_1$, and both with $k=b_1$, $P = \sigma-x_-$ and
$Q = -(\sigma-x_+)$ --- \cref{m:prop:hypIden} reproduces \cref{m:eq:Psihyp} once the
prefactors $\alpha x_+(\sigma-x_-)$ and $-\alpha x_-(\sigma-x_+)$ are carried
through correctly. The factor $(\sigma-x_+)/(\sigma-x_-)$ attached to the second
piece is easily lost, and its loss is invisible at $t=0$; checking $G(1,1,t)=1$
catches it immediately.
\end{remark}
